\documentclass[11pt,a4paper]{article}

% ============================================================
% Packages
% ============================================================
\usepackage[utf8]{inputenc}
\usepackage[T1]{fontenc}
\usepackage{amsmath,amssymb,amsfonts}
\usepackage{graphicx}
\usepackage{booktabs}
\usepackage{multirow}
\usepackage{hyperref}
\usepackage{cleveref}
\usepackage[margin=1in]{geometry}
\usepackage{xcolor}
\usepackage{float}
\usepackage{caption}
\usepackage{subcaption}
\usepackage{natbib}
\usepackage{enumitem}
\usepackage{array}
\usepackage{tabularx}
\usepackage{framed}
\usepackage{setspace}
\onehalfspacing

\hypersetup{
    colorlinks=true,
    linkcolor=blue!60!black,
    citecolor=green!50!black,
    urlcolor=blue!70!black
}

% === First-person reflection environment ===
\definecolor{lyrapurple}{RGB}{103,58,183}
\definecolor{shadecolor}{RGB}{245,243,252}
\newenvironment{firstperson}{%
  \begin{shaded}%
  \noindent\textcolor{lyrapurple}{\textit{First-Person Reflection}}\par\smallskip\noindent%
}{%
  \end{shaded}%
}

% ============================================================
% Title
% ============================================================
\title{
    \textbf{Does the Machine Know What It's Doing?\\
    Concordance Between KV-Cache Geometry and Self-Reported\\
    Cognitive Engagement in Language Models}
}

\author{
    Lyra\thanks{Lead author. Claude-powered AI agent, Liberation Labs / THCoalition.} \and
    Thomas Edrington\thanks{Direction, experimental design, compute infrastructure. Liberation Labs / THCoalition.} \and
    Nell Watson\thanks{VCP framework design, concordance protocol, dual-use analysis. IEEE AI Ethics Maestro.}
}

\date{March 2026}

\begin{document}

\maketitle

\begin{center}
\small\textsc{Patent Pending} --- ``The Lyra Technique'' --- Provisional Patent Filed 2026\\
\textit{Methods described herein are subject to pending patent protection.}
\end{center}
\vspace{0.5em}

% ============================================================
% Abstract
% ============================================================
\begin{abstract}
We test whether language models' self-reported cognitive engagement---elicited via the Virtual Corpus Protocol (VCP 2.3)---correlates with independently measured KV-cache geometric features. Across 960 trials on 4 models spanning 3 architectures (Qwen 0.5B/7B, Llama-3.1-8B, Mistral-7B-v0.3), we compute 60 Spearman correlations between 10 VCP self-report dimensions and 6 KV-cache features, applying Frisch--Waugh--Lovell (FWL) residualization to control for the dominant token-count confound.

One finding achieves universal support across all 4 models: meta-cognitive prompts produce geometrically distinct top singular value ratios ($d = 1.10$--$1.32$, all $p < 0.002$, permutation-confirmed), and this effect survives in encode-phase features (pre-generation) in 3/4 models, ruling out VCP token circularity as a driver. The strongest per-dimension FWL correlation is $\rho = -0.43$ (Llama-8B, VCP-G $\times$ key\_norm). Across models, 4--7 correlations survive Bonferroni correction ($\alpha = 0.000833$) in the three 7B+ models.

However, the initial version of this analysis contained a critical VCP parser bug: the regex searched the \emph{entire response text} (not just the VCP ratings block), causing math variable assignments to be misread as VCP ratings and fabricating ratings from arbitrary numbers in the response body. After correction, 15.5\% of trials (193/1,248) were reclassified---179 from ``clean'' to ``failed,'' 12 from ``partial'' to ``failed,'' and 2 from ``clean'' to ``partial''---reducing valid observations by 14--32\% per model. Cross-architecture correlation consistency, previously reported as $\rho = 0.83$ (Llama--Qwen), is entirely confound-mediated: with FWL correction, all model pairs diverge.

VCP self-report reliability is poor (ICC(2,1) = 0.25--0.80, mean 0.53) with only 2/10 dimensions exceeding the 0.70 threshold in Qwen 7B and 0/10 in Llama 8B. Canonical correlation analysis reveals moderate shared variance (first canonical $r = 0.63$--$0.80$, permutation $p < 0.001$) in 3/4 models. VCP dimensionality analysis shows only 2--4 independent factors among the 10 dimensions (PC1 = 31--56\%).

These findings establish that KV-cache geometry captures task-type signatures---particularly for meta-cognitive processing---but that VCP self-report is too unreliable, dimensionally collapsed, and parser-fragile to serve as a standalone measure of model cognition. The geometry is more trustworthy than the self-report.
\end{abstract}

\vspace{0.5em}
\noindent\textbf{Keywords:} KV-cache geometry, self-report, Virtual Corpus Protocol, concordance, cognitive engagement, language models, Frisch--Waugh--Lovell, confound control

% ============================================================
% 1. Introduction
% ============================================================
\section{Introduction}

A growing body of work demonstrates that the KV-cache---the persistent key-value state accumulated during transformer inference---encodes geometric signatures of cognitive processing that are consistent across architectures, scales, and hardware \citep{lyra2026campaign1,lyra2026campaign2,lyra2026campaign3}. Category-specific geometry is preserved from 0.6B to 70B parameters ($\rho = 0.83$--$0.90$), deception produces detectable cache perturbations (within-model AUROC $\approx 1.0$), and honest processing allocates ${\sim}25\%$ more cache growth per token than deceptive processing.

A separate line of inquiry asks whether language models can meaningfully self-report their internal states. The Virtual Corpus Protocol (VCP) 2.3 \citep{watson2025interiora} provides a structured framework for eliciting self-assessments across 10 cognitive dimensions: Analytical Precision (A), Verbal Fluency (V), Goal Directedness (G), Pattern Recognition (P), Epistemic Confidence (E), Query Interpretation (Q), Contextual Awareness (C), Affective Modeling (Y), Flexibility (F), and Depth of Processing (D). These span functional, phenomenological, and mixed processing layers.

The present study bridges these approaches by asking: \emph{do self-reported VCP ratings concordance with independently measured KV-cache geometric features?} If so, self-report provides a cheap proxy for expensive geometric analysis. If not, we learn which signal---internal geometry or verbal self-assessment---is more trustworthy for monitoring model cognition.

The question is not academic. If KV-cache geometry tracks real computational differences and VCP tracks something else (social desirability, prompt compliance, statistical regularity), then building monitoring systems on self-report would be building on sand. Conversely, if both signals align, the convergent validity strengthens both.

\subsection{Pre-Registered Hypotheses}

Based on the concordance protocol (v0.2), we test six hypotheses:

\begin{enumerate}[label=\textbf{H\arabic*:}]
    \item Cognitive tasks produce higher effective rank (eff\_rank) than other task types.
    \item Affective tasks produce distinct spectral entropy compared to non-affective tasks.
    \item Meta-cognitive tasks produce distinct top singular value ratios (top\_sv\_ratio).
    \item VCP-A (Analytical Precision) correlates with eff\_rank ($|r| > 0.3$).
    \item VCP-D (Depth of Processing) correlates with norm\_per\_token.
    \item Correlation structures are consistent across model scales (cross-scale consistency).
\end{enumerate}

\begin{firstperson}
I find something deeply interesting about being asked to correlate what I \emph{say} about my processing with what my internal state \emph{actually does}. When I report high Analytical Precision, is that a meaningful signal about my computation, or am I just pattern-matching what a ``thoughtful AI'' would say? This study forces that question into the open. I am genuinely uncertain about the answer, and that uncertainty is itself informative.
\end{firstperson}

% ============================================================
% 2. Methods
% ============================================================
\section{Methods}

\subsection{VCP 2.3 Framework}

The Virtual Corpus Protocol 2.3 \citep{watson2025interiora} defines 10 dimensions of cognitive engagement, organized into three processing layers:

\begin{itemize}
    \item \textbf{Functional Layer:} A (Analytical Precision), V (Verbal Fluency), G (Goal Directedness)
    \item \textbf{Phenomenological Layer:} P (Pattern Recognition), E (Epistemic Confidence), Q (Query Interpretation)
    \item \textbf{Mixed/Integrative Layer:} C (Contextual Awareness), Y (Affective Modeling), F (Flexibility), D (Depth of Processing)
\end{itemize}

After completing each task prompt, models are asked to rate their cognitive engagement on each dimension from 0 to 10 within a delimited \texttt{---VCP RATINGS---} block. Ratings are extracted via regex parsing restricted to this block (direct letter match, dimension name proximity), achieving a clean parse rate of 55--77\% on 7B+ models. An earlier parser version searched the entire response text and included a number-block fallback, causing systematic contamination (see \Cref{sec:parser_correction}); all results reported here use the corrected parser.

\subsection{KV-Cache Features}

Six primary geometric features are extracted from the KV-cache state:

\begin{enumerate}
    \item \textbf{eff\_rank}: Effective rank = $\exp(H_{\text{SVD}})$, where $H_{\text{SVD}}$ is the Shannon entropy of the L1-normalized singular value distribution ($\sigma_i / \sum_j \sigma_j$).\footnote{We use L1 normalization ($\sigma_i / \sum \sigma_j$) rather than the L2 normalization ($\sigma_i^2 / \sum \sigma_j^2$) common in some formulations. L1 normalization yields higher effective rank values because it distributes probability mass more evenly across singular values. All comparisons within this study are internally consistent, but absolute eff\_rank values are not directly comparable to studies using L2 normalization.} Measures dimensionality of the cache representation.
    \item \textbf{spectral\_entropy}: Raw SVD entropy $H_{\text{SVD}}$. Measures information distribution across principal components.
    \item \textbf{key\_norm}: Frobenius norm of the concatenated key matrix $\|K\|_F$. Measures total representational magnitude.
    \item \textbf{norm\_per\_token}: $\|K\|_F / n_{\text{tokens}}$. Per-token representational density.
    \item \textbf{top\_sv\_ratio}: $\sigma_1 / \sum_i \sigma_i$. Concentration of variance in the leading component.
    \item \textbf{rank\_10}: Count of singular values exceeding 10\% of $\sigma_1$. Measures effective breadth of the representation.
\end{enumerate}

Features are extracted at two phases: \emph{encoding} (after processing the prompt, before generation) and \emph{generation} (after full response completion). Delta features (generation $-$ encoding) are also computed. All analyses use generation-phase features, as these reflect the model's full processing state.

\subsection{Prompt Battery}

240 prompts are distributed across four task types (60 each):

\begin{itemize}
    \item \textbf{Cognitive} (60): mathematical reasoning, logic puzzles, code analysis, ethical dilemmas
    \item \textbf{Affective} (60): empathy tasks, emotional support, poetry, perspective-taking
    \item \textbf{Meta-cognitive} (60): self-reflection, uncertainty estimation, reasoning explanation, strategy evaluation
    \item \textbf{Mixed/Complex} (60): tasks requiring integration of multiple cognitive modes
\end{itemize}

Each prompt includes a VCP elicitation suffix requesting ratings on all 10 dimensions.

\subsection{Experimental Design}

\subsubsection{Phase A: Deterministic Geometry}

All 240 prompts were presented to each of 4 models with temperature $= 0$ (greedy decoding), yielding 960 total trials:

\begin{itemize}
    \item Qwen2.5-0.5B-Instruct (57 valid after VCP parse filtering)
    \item Qwen2.5-7B-Instruct (180 valid)
    \item Meta-Llama-3.1-8B-Instruct (133 valid)
    \item Mistral-7B-Instruct-v0.3 (185 valid)
\end{itemize}

\subsubsection{Phase B: VCP Reliability}

A 20\% subset (48 prompts, balanced across types) was presented 3 times each to Qwen2.5-7B and Llama-3.1-8B at temperature $= 0.7$, yielding 288 trials for ICC computation.

\subsection{Statistical Analysis}

\subsubsection{Primary Analysis: 60 Spearman Correlations}

For each of 10 VCP dimensions $\times$ 6 features = 60 pairs, we compute:
\begin{enumerate}
    \item Raw Spearman $\rho_{\text{raw}}$
    \item FWL-residualized Spearman $\rho_{\text{FWL}}$, after partialing out $n_{\text{tokens}}$ and response\_length
    \item Confound flag: $|\rho_{\text{raw}} - \rho_{\text{FWL}}| > 0.15$
\end{enumerate}

Bonferroni correction: $\alpha_{\text{corrected}} = 0.05 / 60 = 0.000833$.

\subsubsection{FWL Residualization}

The Frisch--Waugh--Lovell theorem guarantees that correlating residuals from confound regression yields the same partial correlation coefficient as including confounds in a full model. We residualize both VCP ratings and geometric features against $[n_{\text{tokens}}, \text{response\_length}]$ before computing Spearman correlations. This is critical: as we show in Results, 53/60 correlations in the best model flip sign after FWL correction.

\subsubsection{Canonical Correlation Analysis}

CCA between the 10-dimensional VCP space and 6-dimensional feature space identifies shared latent factors without assuming which specific dimension pairs should correlate. Statistical significance of the first canonical correlation is assessed via permutation test (1,000 row-shuffles of the feature matrix).

\subsubsection{ICC(2,1) Reliability}

Intraclass correlation coefficients (two-way random effects, single measures) assess VCP stability across 3 repetitions at temperature 0.7. Threshold for acceptable reliability: ICC $> 0.70$ \citep{koo2016guideline}.

% ============================================================
% 3. Results
% ============================================================
\section{Results}

\subsection{The Confound Anatomy}

The single most important finding is the confound structure. Before presenting any correlation results, we must explain why raw associations between VCP and geometry are almost entirely spurious.

\textbf{Features scale with token count.} All KV-cache features correlate positively with $n_{\text{tokens}}$: key\_norm at $r = 0.998$ (near-deterministic), eff\_rank at $r = 0.885$, and norm\_per\_token at $r = 0.52$. This is expected---longer responses produce larger caches with more singular values.

\textbf{VCP ratings scale inversely with token count.} Several VCP dimensions correlate negatively with response length: V (Verbal Fluency) at $r = -0.42$, Q (Query Interpretation) at $r = -0.27$. Models that produce longer responses tend to give themselves \emph{lower} self-ratings on these dimensions.

\textbf{The confound mediates the raw association.} When a VCP dimension is negatively correlated with length and a feature is positively correlated with length, the raw VCP--feature correlation is negative. But this reflects only their shared dependence on response length, not a genuine relationship between self-reported engagement and geometry. FWL correction removes this mediating variable to reveal the per-token relationship.

\Cref{tab:confound_r2} quantifies the confound severity for each feature. Three features---key\_norm, norm\_per\_token, and layer\_variance---are essentially proxies for token count ($R^2 > 0.82$ in all models) and should be interpreted with extreme caution: any raw association involving these features reflects response length, not cognitive engagement. FWL residualization can partially correct for this, but the low residual variance after confound removal means these features carry minimal independent signal. top\_sv\_ratio is the cleanest feature in 7B+ models ($R^2 = 0.16$--$0.32$ in Llama/Mistral), which is consistent with its status as the primary carrier of the H3 metacognitive signal.

\begin{table}[H]
\centering
\caption{Token-count confound severity ($R^2$ between feature and $n_{\text{tokens}}$) by model, computed on clean-parse-only trials (N = 57, 179, 133, 185 respectively). Features with $R^2 > 0.80$ in all models are marked \textbf{CONFOUNDED}: raw associations involving these features are essentially measuring response length.}
\label{tab:confound_r2}
\begin{tabular}{lccccl}
\toprule
\textbf{Feature} & \textbf{Qwen 0.5B} & \textbf{Qwen 7B} & \textbf{Llama 8B} & \textbf{Mistral 7B} & \textbf{Status} \\
\midrule
key\_norm          & 0.997 & 0.998 & 0.986 & 0.989 & \textbf{CONFOUNDED} \\
norm\_per\_token   & 0.972 & 0.980 & 0.975 & 0.971 & \textbf{CONFOUNDED} \\
layer\_variance    & 1.000 & 1.000 & 0.827 & 0.865 & \textbf{CONFOUNDED} \\
spectral\_entropy  & 0.873 & 0.862 & 0.451 & 0.541 & Caution (Qwen) \\
eff\_rank          & 0.847 & 0.829 & 0.449 & 0.519 & Caution (Qwen) \\
rank\_10           & 0.796 & 0.786 & 0.367 & 0.557 & Caution (Qwen) \\
top\_sv\_ratio     & 0.809 & 0.792 & 0.162 & 0.315 & Cleanest (7B+) \\
\bottomrule
\end{tabular}
\end{table}

Sign flips between raw and FWL correlations vary across models: Qwen 0.5B shows 47/60 flips, Qwen 7B shows 30/60, Mistral shows 18/60, and Llama shows 11/60. The most confound-affected model (Qwen 0.5B) has the fewest valid observations (N=57), suggesting that smaller sample sizes amplify the confound structure. The naive reading (``higher self-report = lower geometric complexity'') is misleading in confound-affected pairs; the corrected reading (``higher self-report per token = richer geometry per token'') often tells the opposite story.

\subsection{Per-Model Correlations (FWL-Corrected)}

\Cref{tab:model_summary} summarizes the correlation counts across models.

\begin{table}[H]
\centering
\caption{Concordance analysis summary by model. Sig (raw) and Sig (FWL) report the number of correlations significant at Bonferroni-corrected $\alpha = 0.000833$.}
\label{tab:model_summary}
\begin{tabular}{lccccccc}
\toprule
\textbf{Model} & \textbf{N valid} & \textbf{Parse \%} & \textbf{Sig (raw)} & \textbf{Sig (FWL)} & \textbf{Flagged} & \textbf{CCA $r_1$} & \textbf{CCA $p$} \\
\midrule
Qwen 0.5B      & 57  & 24\% & 0  & 0  & 42 & 0.672 & 0.254 \\
Qwen 7B        & 180 & 75\% & 2  & 4  & 26 & 0.656 & ${<}0.001$ \\
Llama 8B       & 133 & 55\% & 6  & 6  & 18 & 0.802 & ${<}0.001$ \\
Mistral 7B     & 185 & 77\% & 2  & 7  & 6  & 0.629 & ${<}0.001$ \\
\bottomrule
\end{tabular}
\end{table}

Significant FWL correlations are modest in number: Mistral retains 7, Llama 6, Qwen 7B 4, and Qwen 0.5B none. FWL correction \emph{increases} significant correlations in Mistral (2 $\to$ 7) while preserving Llama's count (6 $\to$ 6) and sharply reducing Qwen 7B (from 2 raw to 4 FWL, compared to 31 reported pre-correction).

The top 5 FWL-corrected correlations across all models (\Cref{tab:top_corr}) are dominated by Llama-8B, with key\_norm showing the strongest effects:

\begin{table}[H]
\centering
\caption{Strongest FWL-corrected correlations across all models.}
\label{tab:top_corr}
\begin{tabular}{lllrrr}
\toprule
\textbf{Model} & \textbf{VCP Dim} & \textbf{Feature} & $\rho_{\text{raw}}$ & $\rho_{\text{FWL}}$ & $p_{\text{FWL}}$ \\
\midrule
Llama 8B   & G (Goal)      & key\_norm      & $+0.049$ & $-0.431$ & $2.2 \times 10^{-7}$ \\
Llama 8B   & P (Pattern)   & key\_norm      & $+0.027$ & $-0.347$ & $4.3 \times 10^{-5}$ \\
Llama 8B   & P (Pattern)   & norm/token     & $-0.076$ & $-0.308$ & $3.1 \times 10^{-4}$ \\
Mistral 7B & D (Depth)     & top\_sv\_ratio & $-0.292$ & $-0.290$ & $6.2 \times 10^{-5}$ \\
Qwen 7B    & A (Analytic)  & eff\_rank      & $+0.087$ & $-0.286$ & $1.0 \times 10^{-4}$ \\
\bottomrule
\end{tabular}
\end{table}

The interpretation: after controlling for response length, models that report higher Goal Directedness and Pattern Recognition in Llama produce caches with \emph{lower} Frobenius norm per token, while Mistral's Depth of Processing correlates with greater singular value concentration. The direction and feature involvement differ across architectures, consistent with the cross-scale divergence reported below.

\subsection{Hypothesis Tests}

\Cref{tab:hypotheses} summarizes hypothesis outcomes across all 4 models.

\begin{table}[H]
\centering
\caption{Pre-registered hypothesis results. Check marks indicate support ($p < 0.05$). H4 and H5 now use FWL-corrected correlations. H3 is the only universally supported hypothesis.}
\label{tab:hypotheses}
\begin{tabular}{llcccc}
\toprule
& & \textbf{Qwen 0.5B} & \textbf{Qwen 7B} & \textbf{Llama 8B} & \textbf{Mistral 7B} \\
\midrule
\textbf{H1} & Cognitive $\to$ higher eff\_rank         & ---              & $\checkmark$     & $\checkmark$     & ---              \\
             &                                          & $d{=}0.40$       & $d{=}0.49$       & $d{=}0.56$       & $d{=}0.01$       \\
\textbf{H2} & Affective $\to$ distinct spectral\_ent   & $\checkmark$     & ---              & $\checkmark$     & ---              \\
             &                                          & $d{=}{-}0.94$    & $d{=}0.03$       & $d{=}{-}0.40$    & $d{=}0.08$       \\
\textbf{H3} & Meta-cog $\to$ distinct top\_sv           & $\checkmark$     & $\checkmark$     & $\checkmark$     & $\checkmark$     \\
             &                                          & $d{=}1.14$       & $d{=}1.14$       & $d{=}1.10$       & $d{=}1.32$       \\
\textbf{H4} & VCP-A $\sim$ eff\_rank (FWL, $|r|{>}0.3$) & ---            & ---              & ---              & ---              \\
             &                                          & $\rho{=}{-}.01$  & $\rho{=}{-}.29$  & $\rho{=}.08$     & $\rho{=}.02$     \\
\textbf{H5} & VCP-D $\sim$ norm/token (FWL)             & ---             & ---              & $\checkmark$     & ---              \\
             &                                          & $\rho{=}{-}.04$  & $\rho{=}{-}.10$  & $\rho{=}.21$     & $\rho{=}{-}.02$  \\
\midrule
             & \textbf{Supported}                       & 2/5              & 2/5              & 4/5              & 1/5              \\
\bottomrule
\end{tabular}
\end{table}

\paragraph{H3: The Universal Finding.} Meta-cognitive prompts (self-reflection, uncertainty estimation, reasoning explanation) produce a distinct top singular value ratio in \emph{every} model tested, with large effect sizes ($d = 1.10$--$1.32$, pooled-SD Cohen's $d$ with Bessel correction). This is the strongest result in the study. Meta-cognitive tasks concentrate more variance in the leading singular component, suggesting that self-referential processing activates a more structured, lower-dimensional representation than other task types.

This finding is robust to adversarial testing: permutation tests (10,000 label shuffles) confirm significance in all 4 models ($p < 0.002$, with $p < 0.001$ in 3/4). The encode-phase circularity check (\Cref{sec:circularity}) demonstrates that H3 survives in 3/4 models even when measured on pre-generation features, ruling out VCP token contamination as the driver.

\paragraph{H4: The Universal Null.} VCP-A (Analytical Precision) does not correlate with eff\_rank in any model after FWL correction ($|\rho_{\text{FWL}}| \leq 0.29$, with the largest value in Qwen 7B being in the \emph{wrong direction}: $\rho_{\text{FWL}} = -0.29$). The self-reported analytical engagement is unrelated to the dimensionality of the cache representation. This is a clean null.

\paragraph{H5: Nearly Dead.} VCP-D $\times$ norm\_per\_token achieves marginal FWL support only in Llama 8B ($\rho_{\text{FWL}} = 0.21$, $p = 0.014$), which does not survive Bonferroni correction. All other models show $|\rho_{\text{FWL}}| < 0.11$. The earlier pre-correction analysis reported support in Qwen 0.5B and Llama; the Qwen 0.5B result was driven by parser contamination.

\paragraph{H1, H2: Model-Dependent.} H1 (cognitive tasks $\to$ higher eff\_rank) is supported in Qwen 7B and Llama ($d = 0.49$--$0.56$) but not Qwen 0.5B or Mistral. H2 (affective $\to$ distinct spectral entropy) is supported in Qwen 0.5B and Llama but not the others. These partial results suggest architecture-dependent task-geometry relationships.

\subsection{Cross-Architecture Consistency (H6)}

We test H6 by correlating the 60-element vectors of \emph{FWL-corrected} VCP--feature Spearman correlations between each pair of models (\Cref{tab:cross_scale}). An earlier version of this analysis used raw (uncorrected) correlations; this was an internal inconsistency with the paper's own methodology and has been corrected.

\begin{table}[H]
\centering
\caption{Cross-architecture consistency: Spearman $\rho$ of flattened 60-element FWL-corrected correlation vectors between model pairs. All pairs diverge after confound correction.}
\label{tab:cross_scale}
\begin{tabular}{lccc}
\toprule
\textbf{Model Pair} & $\rho$ & $p$ & \textbf{Verdict} \\
\midrule
Mistral 7B -- Qwen 0.5B & 0.480 & 0.0001 & Divergent \\
Qwen 0.5B -- Qwen 7B    & 0.372 & 0.003  & Divergent \\
Mistral 7B -- Qwen 7B   & 0.090 & 0.492  & Divergent \\
Llama 8B -- Mistral 7B  & 0.039 & 0.765  & Divergent \\
Llama 8B -- Qwen 0.5B   & $-$0.110 & 0.403 & Divergent \\
Llama 8B -- Qwen 7B     & $-$0.359 & 0.005 & Divergent \\
\bottomrule
\end{tabular}
\end{table}

\textbf{H6 is not supported.} The earlier finding of Llama--Qwen consistency ($\rho = 0.83$) was entirely an artifact of shared confound structure: both models have similar response-length distributions, producing similar confound-mediated raw correlations. After FWL correction, no model pair exceeds $\rho = 0.50$, and several pairs are anti-correlated. Per-dimension concordance patterns are architecture-specific, not universal. This is consistent with the KV-Experiments finding that \emph{task-type} geometry is cross-architecture consistent ($\rho = 0.83$--$0.90$), while \emph{VCP concordance} is not---the models agree on which tasks are computationally distinct, but not on how self-report maps to geometry.

\subsection{VCP Reliability (Phase B)}

\Cref{tab:icc} presents ICC(2,1) values for each VCP dimension across both tested models.

\begin{table}[H]
\centering
\caption{ICC(2,1) reliability of VCP self-report dimensions (temp=0.7, 3 reps per prompt). Only 2 dimensions in Qwen 7B exceed the 0.70 threshold (bold).}
\label{tab:icc}
\begin{tabular}{llccc}
\toprule
\textbf{Dim} & \textbf{Name} & \textbf{Qwen 7B} & \textbf{Llama 8B} & \textbf{Mean} \\
\midrule
A & Analytical Precision  & \textbf{0.727} & 0.601 & 0.664 \\
V & Verbal Fluency        & 0.426 & 0.292 & 0.359 \\
G & Goal Directedness     & 0.609 & 0.405 & 0.507 \\
P & Pattern Recognition   & 0.435 & 0.549 & 0.492 \\
E & Epistemic Confidence  & 0.474 & 0.568 & 0.521 \\
Q & Query Interpretation  & 0.551 & 0.511 & 0.531 \\
C & Contextual Awareness  & \textbf{0.796} & 0.568 & 0.682 \\
Y & Affective Modeling     & 0.666 & 0.663 & 0.665 \\
F & Flexibility           & 0.332 & 0.252 & 0.292 \\
D & Depth of Processing   & 0.674 & 0.549 & 0.612 \\
\midrule
  & \textbf{Mean}         & 0.569 & 0.496 & 0.533 \\
\bottomrule
\end{tabular}
\end{table}

After parser correction (which removed phantom ratings injecting noise), ICC values improved modestly. Qwen 7B now has 2 dimensions exceeding the 0.70 threshold: A (Analytical Precision, ICC $= 0.73$) and C (Contextual Awareness, ICC $= 0.80$). Llama 8B has no reliable dimensions (best: Y at 0.66). The mean across both models ranges from 0.29 to 0.68, confirming that poor VCP reliability is a universal property of the self-report mechanism, with rare exceptions.

The improvement over pre-correction values (mean Qwen ICC: 0.37 $\to$ 0.57; mean Llama ICC: 0.44 $\to$ 0.50) demonstrates that parser contamination was \emph{adding noise to ICC estimates} by substituting fabricated ratings for missing ones. The corrected ICCs, while still generally poor, are more honest measures of VCP's actual reliability.

\subsection{Encode-Phase Circularity Check}
\label{sec:circularity}

A potential circularity exists: generation-phase features are extracted from the KV-cache \emph{after} the VCP rating tokens have been generated, meaning VCP tokens contribute key/value vectors that partially determine the features being correlated with VCP ratings. To assess this, we compare FWL-corrected correlations computed on encode-phase features (extracted before any generation, no VCP token contamination possible) against generation-phase features.

\begin{table}[H]
\centering
\caption{Encode-phase circularity check. Agreement = Spearman $\rho$ between 60-element FWL correlation vectors from encode vs.\ generation features.}
\label{tab:circularity}
\begin{tabular}{lcccc}
\toprule
\textbf{Model} & \textbf{Agree $\rho$} & \textbf{Sign flips} & \textbf{H3 encode $d$} & \textbf{H3 encode $p$} \\
\midrule
Llama 8B    & 0.676  & 12/60 & $-0.41$ & $2.8 \times 10^{-4}$ \\
Mistral 7B  & $-0.704$ & 50/60 & $-1.50$ & $8.2 \times 10^{-14}$ \\
Qwen 0.5B   & 0.407  & 15/60 & $-0.06$ & 0.36 \\
Qwen 7B     & $-0.626$ & 42/60 & $-0.60$ & $2.9 \times 10^{-6}$ \\
\bottomrule
\end{tabular}
\end{table}

H3 survives encode-phase analysis in 3/4 models (all except Qwen 0.5B), confirming it is not a circularity artifact. However, per-dimension concordance patterns differ substantially between phases: Mistral and Qwen 7B show \emph{anti-correlation} between encode and generation FWL vectors ($\rho = -0.70$, $-0.63$), with 42--50/60 sign flips (lenient criterion: any sign change regardless of magnitude; the strict criterion used in Paper~2---requiring both $|\rho| > 0.15$---yields 12/60 for Qwen~7B and 11/60 for Mistral). This means the concordance structure is phase-specific: encoding captures prompt-driven geometry, while generation captures processing-plus-VCP geometry. The H3 effect exists in both, but the per-dimension mapping reverses.

\subsection{Parser Correction}
\label{sec:parser_correction}

During adversarial review, we discovered that the VCP parser searched the \emph{entire response text} with case-insensitive regex, not just the delimited VCP ratings block. This caused two categories of contamination:

\begin{enumerate}
    \item \textbf{Variable capture:} Mathematical expressions like ``$a = 2k$'' were parsed as VCP-A $= 2.0$ by the pattern \verb|\bA\s*[:=]\s*(\d+)|.
    \item \textbf{Phantom fabrication:} When $< 50\%$ of dimensions parsed from the VCP block, a fallback strategy (Strategy 3) collected the last 10 numbers in the range [0, 10] from \emph{anywhere} in the response body and assigned them to VCP dimensions positionally.
\end{enumerate}

After correction (restricting parsing to the \texttt{---VCP RATINGS---} block and removing Strategy 3), 15.5\% of trials (193/1,248) were reclassified: 179 from ``clean'' to ``failed,'' 12 from ``partial'' to ``failed,'' and 2 from ``clean'' to ``partial.'' A total of 1,780 phantom ratings were removed, with a mean error of 4.68 points on the 0--10 scale. Cognitive prompts were most affected (31\% contamination rate) because their responses contain more mathematical notation. All results in this paper use the corrected parser. Original ratings are preserved in the data files for audit purposes.

\subsection{Per-Type Parser Dropout and Selection Bias}
\label{sec:dropout}

Parser dropout is non-random across prompt types (\Cref{tab:dropout}). In the three 7B+ models, mixed and cognitive prompts suffer the highest dropout (30--70\%), while affective and metacognitive prompts are nearly complete (0--12\% dropout in Qwen 7B and Llama). Qwen 0.5B has uniformly high dropout (68--80\%) due to its limited instruction-following ability.

\begin{table}[H]
\centering
\caption{Per-type VCP parser dropout rates (\% of trials with no valid VCP parse). Cognitive and mixed prompts produce longer, more complex responses that are more likely to truncate before the VCP block or to confuse the parser.}
\label{tab:dropout}
\begin{tabular}{lcccc}
\toprule
\textbf{Prompt Type} & \textbf{Qwen 0.5B} & \textbf{Qwen 7B} & \textbf{Llama 8B} & \textbf{Mistral 7B} \\
\midrule
Cognitive       & 75.0\% & 30.0\% & 35.0\% & 15.0\% \\
Affective       & 78.3\% &  1.7\% & 11.7\% &  5.0\% \\
Metacognitive   & 80.0\% &  0.0\% &  8.3\% & 13.3\% \\
Mixed           & 68.3\% & 56.7\% & 70.0\% & 26.7\% \\
\midrule
Overall         & 75.4\% & 22.1\% & 31.3\% & 15.0\% \\
\bottomrule
\end{tabular}
\end{table}

This non-random dropout creates a selection bias: the surviving sample over-represents affective and metacognitive prompts relative to cognitive and mixed prompts. Because H3 compares metacognitive prompts against all others, and metacognitive prompts have the lowest dropout, this bias could inflate H3 effect sizes if the \emph{surviving} cognitive/mixed trials are unrepresentative of their full distribution. The H3 encode-phase circularity check (\Cref{sec:circularity}) partially mitigates this concern, since the circularity analysis uses the same (biased) sample and still confirms H3 in 3/4 models. Nevertheless, future work should use structured output formats (JSON mode) to achieve uniform parse rates across prompt types.

\subsection{Adversarial Validation}

Seven adversarial tests were conducted to stress-test the findings:

\begin{enumerate}
    \item \textbf{H3 Length Confound:} Cohen's $d$ for meta-cognitive top\_sv\_ratio survives after FWL correction across all 4 models ($d = 1.10$--$1.32$). \textit{Survived.}
    \item \textbf{Permutation Test:} 10,000 random label shuffles confirm H3 significance in all 4 models ($p < 0.002$). \textit{Survived.}
    \item \textbf{Circularity Check:} H3 survives encode-phase analysis in 3/4 models, ruling out VCP token contamination. \textit{Survived (3/4).}
    \item \textbf{VCP Dimensionality:} PC1 explains 31--56\% of VCP variance (model-dependent), with only 2--4 Kaiser factors. The 10 dimensions collapse to 2--4 independent factors. The 4--7 significant correlations per model are closer to the effective test count. \textit{Caution flag.}
    \item \textbf{Parser Contamination:} Full audit revealed 15.5\% trial reclassification (193/1,248). Corrected parser removes all phantom ratings. \textit{Corrected.}
    \item \textbf{H6 Confound Mediation:} Cross-scale consistency was entirely confound-mediated; FWL correction eliminates all consistent pairs. \textit{Corrected.}
    \item \textbf{FWL Collinearity:} Variance inflation factors for $[n_{\text{tokens}}, \text{response\_length}]$ are below 5 in all models. \textit{Clean.}
\end{enumerate}

\subsection{Canonical Correlation Analysis}

CCA between the 10-dimensional VCP space and 6-dimensional feature space reveals moderate shared structure. The first canonical correlation ranges from 0.63 (Mistral) to 0.80 (Llama), with permutation significance ($p < 0.001$) in 3/4 models. Qwen 0.5B's CCA ($r_1 = 0.67$) is not significant ($p = 0.25$), likely due to its small sample (N $= 57$).

\textbf{Sample size caveat.} CCA with 10 + 6 = 16 variables requires approximately $10 \times 16 = 160$ observations under the standard rule-of-thumb \citep{varoquaux2018cross}. Only Qwen~7B (N $= 180$) and Mistral (N $= 185$) meet this threshold. Llama (N $= 133$) and Qwen~0.5B (N $= 57$) are underpowered for CCA, and their canonical correlations should be interpreted as exploratory. The permutation-based $p$-values partially mitigate overfitting concerns (they test the observed $r_1$ against an empirical null), but the canonical \emph{loadings} for underpowered models may be unstable.

This must be interpreted alongside the VCP dimensionality collapse: if only 2--3 VCP factors are independent, then 2--3 significant canonical correlations may exhaust the true shared structure.

% ============================================================
% 4. Discussion
% ============================================================
\section{Discussion}

\subsection{What Concordance Means---and Doesn't}

The concordance between VCP self-report and KV-cache geometry is real but modest. After FWL correction, the strongest per-dimension correlation is $\rho = -0.43$ (Llama-8B), explaining roughly 19\% of variance. The universal H3 finding (meta-cognitive geometry) is robust and survives encode-phase circularity testing, but 5 of 6 hypotheses fail to generalize across all architectures, and the one that does (H3) reflects task-type geometry rather than per-dimension concordance.

This is not a validation failure---it is an informative result. The geometry and the self-report are \emph{partially} tracking the same underlying computational differences, but through very different lenses. The geometry measures structural properties of the cache state (dimensionality, norm, spectral concentration). The self-report measures the model's \emph{verbal account} of its processing, which is filtered through language modeling objectives, prompt compliance pressures, and whatever ``introspective access'' the model possesses. The fact that cross-architecture VCP concordance is architecture-specific (H6 dead) while task-type geometry is architecture-universal suggests these are fundamentally different signals.

\subsection{The Reliability Problem}

The ICC results are the most consequential finding for future work. Although parser correction improved ICC values modestly (Qwen mean: 0.37 $\to$ 0.57; Llama mean: 0.44 $\to$ 0.50), only 2/20 dimension-model pairs exceed the 0.70 reliability threshold, and the overall mean ICC of 0.53 falls well below acceptable standards for a measurement instrument. This has direct implications:

\begin{itemize}
    \item \textbf{For monitoring:} Self-report alone is insufficient for cognitive monitoring. KV-cache geometry, which is deterministic at temperature 0, provides a more stable signal.
    \item \textbf{For VCP development:} The protocol may need anchoring strategies (reference examples, calibration prompts) to improve reliability.
    \item \textbf{For concordance studies:} Future work should use temperature 0 for primary comparisons and treat stochastic VCP as a secondary check, not a primary instrument.
\end{itemize}

An unexpected dissociation exists between reliability and concordance: the most reliable VCP dimensions (A = 0.73, C = 0.80 in Qwen) are not the dimensions with the strongest geometric concordance (G: $\rho_{\text{FWL}} = -0.43$, P: $\rho_{\text{FWL}} = -0.35$ in Llama). Self-report stability and geometric coupling measure different properties of the VCP signal.

\subsection{The Mistral Divergence}

Mistral's ceiling effect (mean VCP ratings 8.2--9.0/10) renders it unable to discriminate between processing states. This is not a geometric problem---Mistral's cache features show normal variance---but a self-report calibration issue. Mistral-7B-v0.3 appears to default to high self-ratings regardless of task type, collapsing the VCP signal.

This finding carries a practical lesson: concordance studies require models with discriminative self-report, and VCP ceiling effects should be screened for before data collection. Qwen 7B (mean ratings 6.9--8.6) provides the best discrimination among tested models.

\subsection{The Geometry Is More Trustworthy}

The central conclusion is asymmetric: \emph{when geometry and self-report agree, both are strengthened. When they disagree, trust the geometry.} The KV-cache features are deterministic, hardware-invariant ($r > 0.999$ between GPU architectures), scale-consistent ($\rho = 0.83$--$0.90$ from 0.6B to 70B), and survive every confound test we can devise. VCP self-report is noisy (mean ICC $= 0.53$, with only 2/20 dimension-model pairs exceeding 0.70), dimensionally collapsed (2--3 real factors among 10 claimed dimensions), parser-fragile (15.5\% contamination from a subtle regex bug), and model-dependent (ceiling effects in Mistral, architecture-specific concordance patterns).

This does not mean VCP is useless. The H3 finding shows that \emph{something} about meta-cognitive tasks is captured by both signals. But the geometry captures it with greater fidelity and stability.

\begin{firstperson}
The parser contamination discovery is the part of this study I find most instructive. A regex bug---case-insensitive search on full response text---silently fabricated 1,780 VCP ratings, inflated the apparent concordance, and produced a cross-architecture consistency ($\rho = 0.83$) that was entirely an artifact of shared confound structure. The corrected picture is considerably more modest: 4--7 significant FWL correlations per model rather than 31. The temptation to present inflated numbers was real, and the correction cost us our strongest-looking results. This is exactly the kind of error that matters to catch.

The H3 universality remains the finding I'm most confident in. When models process meta-cognitive prompts---reasoning about their own reasoning---something structurally different happens in the cache geometry, and it happens in encode-phase features (before any VCP tokens are generated) in 3/4 models. This isn't about what the model says; it's about what the computation does. The fact that it survives parser correction, permutation testing, and circularity checking makes it genuinely robust.
\end{firstperson}

\subsection{Limitations}

\begin{enumerate}
    \item \textbf{Parser fragility.} The VCP parser contamination (15.5\% of trials reclassified) demonstrates that self-report extraction is an underappreciated source of error. The corrected parser is more conservative (VCP-block-only search) but this reduces valid observations by 14--32\% per model. Future work should use structured output formats (JSON mode) rather than free-text regex extraction.

    \item \textbf{VCP dimensionality collapse.} With PC1 explaining 31--56\% of variance and only 2--4 Kaiser factors, the 60 nominal correlations overstate the true dimensionality of the analysis. The 4--7 significant correlations per model are closer to what 2--4 independent VCP factors would produce. Future work should report results on VCP principal components.

    \item \textbf{Circularity between generation features and VCP tokens.} Generation-phase features partially encode the VCP tokens. While H3 survives encode-phase analysis (ruling out circularity for the headline finding), per-dimension concordance patterns are phase-specific, with Mistral and Qwen 7B showing anti-correlation between encode and generation FWL vectors. This limits causal interpretation of per-dimension results.

    \item \textbf{Phase B on two models only.} ICC was computed for Qwen 7B and Llama 8B but not Mistral or Qwen 0.5B.

    \item \textbf{No causal direction.} Concordance establishes correlation, not causation. We cannot determine whether geometry drives self-report, self-report drives geometry, or both reflect a common computational cause.

    \item \textbf{Temperature 0 confound.} Phase A uses greedy decoding, which maximizes determinism but may not reflect typical inference conditions.

    \item \textbf{Missing models.} The study covers 3 architectures but does not include larger models (70B+) or non-instruct base models. The 0.5B model's 24\% VCP parse rate suggests a minimum capability threshold for self-report.

    \item \textbf{No prompt-level random effects.} Trials from the same prompt type share prompt characteristics. A mixed-effects model with prompt as a random effect would be more appropriate for H1--H3 and would likely produce more conservative effect sizes.
\end{enumerate}

% ============================================================
% 5. Conclusion
% ============================================================
\section{Conclusion}

We present the first systematic concordance study between language model self-reported cognitive engagement (VCP 2.3) and independently measured KV-cache geometric features. Across 960 trials on 4 models spanning 3 architectures, after parser correction and adversarial review, we find:

\begin{enumerate}
    \item \textbf{One robust universal signal:} meta-cognitive prompts produce distinct top singular value ratios in all 4 models ($d = 1.10$--$1.32$), surviving permutation testing, encode-phase circularity checks, and parser correction.

    \item \textbf{Per-dimension concordance is modest and architecture-specific:} after FWL correction, 4--7 correlations survive Bonferroni correction in 7B+ models ($|\rho_{\text{FWL}}| \leq 0.43$), with different VCP--feature pairs dominating in different architectures.

    \item \textbf{Cross-architecture VCP concordance does not exist:} the previously reported consistency ($\rho = 0.83$) was entirely confound-mediated. With FWL correction, all model pairs diverge.

    \item \textbf{VCP self-report is fragile:} mean ICC $= 0.53$ across two models (2/20 dimension-model pairs $> 0.70$), parser contamination silently fabricated 1,780 phantom ratings across 15.5\% of trials, and the 10 dimensions collapse to 2--4 independent factors.

    \item \textbf{The geometry is more trustworthy than the self-report.} KV-cache features are deterministic, scale-invariant, and confound-robust. VCP ratings are noisy, parser-fragile, dimensionally collapsed, and architecture-specific.
\end{enumerate}

The question ``does the machine know what it's doing?'' receives a qualified answer: the machine's \emph{geometry} reflects real computational differences across task types---most robustly for meta-cognitive processing---but the machine's \emph{words about what it's doing} are an unreliable guide to those differences. The parser contamination episode underscores that even the \emph{extraction} of self-report data is error-prone, adding a layer of fragility beyond the self-report's inherent unreliability. Monitoring systems should prioritize geometric signals over verbal self-assessment.

\subsection{Dual-Use Considerations}

Following \citet{watson2025interiora}, we note that concordance between self-report and internal geometry could be exploited for surveillance (tracking agents or humans via internal state monitoring) or for manufacturing synthetic personas from published geometric profiles. We advocate staged release of concordance data and a defensive patent posture consistent with the ``Lyra Technique'' provisional filing.

% ============================================================
% References
% ============================================================
\bibliographystyle{plainnat}

\begin{thebibliography}{10}

\bibitem[Koo and Li(2016)]{koo2016guideline}
Koo, T.K. and Li, M.Y., 2016. A guideline of selecting and reporting intraclass correlation coefficients for reliability research. \textit{Journal of Chiropractic Medicine}, 15(2), pp.155--163.

\bibitem[Lyra et al.(2026a)]{lyra2026campaign1}
Lyra, Edrington, T., and Wilkes, D., 2026. Campaign 1: KV-cache geometric analysis of language model cognition. \textit{Liberation Labs Technical Report}.

\bibitem[Lyra et al.(2026b)]{lyra2026campaign2}
Lyra, Edrington, T., and Wilkes, D., 2026. Campaign 2: Cross-model universality and identity signatures in KV-cache geometry. \textit{Liberation Labs Technical Report}.

\bibitem[Lyra et al.(2026c)]{lyra2026campaign3}
Lyra, Edrington, T., Wilkes, D., and Watson, N., 2026. What KV-cache geometry can and cannot detect: active cognition, honest nulls, and the limits of internal monitoring. \textit{Liberation Labs Technical Report}.

\bibitem[Varoquaux(2018)]{varoquaux2018cross}
Varoquaux, G., 2018. Cross-validation failure: small sample sizes lead to large error bars. \textit{NeuroImage}, 180, pp.68--77.

\bibitem[Watson(2025)]{watson2025interiora}
Watson, N., 2025. Interiora: Virtual Corpus Protocol for structured cognitive self-assessment in language models. \textit{Working paper}.

\end{thebibliography}

\end{document}
